\documentclass{article}
\usepackage[T1]{fontenc}
\usepackage{lmodern}
\usepackage{graphicx}
\usepackage{amsmath,amssymb}
\usepackage[a4paper,margin=2cm]{geometry}
\usepackage[absolute,overlay]{textpos}
\setlength{\TPHorizModule}{1mm}
\setlength{\TPVertModule}{1mm}
\usepackage[indonesian]{babel}
\usepackage{hyperref}
\setlength{\emergencystretch}{2em}
% unit: Halaman 44

\begin{document}

% === Halaman 44 (llm) ===

\section*{2. Serangan aktif (active attack)}

\begin{itemize}
    \item penyerang mengintervensi komunikasi dan ikut mempengaruhi sistem untuk keuntungan dirinya
    \item penyerang dapat mengubah aliran pesan seperti:
    \begin{itemize}
        \item menghapus sebagian cipherteks,
        \item mengubah cipherteks,
        \item menyisipkan potongan cipherteks palsu,
        \item me-replay pesan lama,
        \item mengubah informasi yang tersimpan, dsb
    \end{itemize}
    \item Contoh: \textit{man-in-the-middle attack}
\end{itemize}

\end{document}
